\section{A jointly entire kernel for all zero-block decorations}
\label{zbg:section}

The one-free-order zero-block generator in the incoming harmonic-parity
report admits a uniform extension to any number of independently moving
orders.  The extension is not obtained by interpolating finite sums from
integer endpoint differences.  A divided-difference identity proves the
required analytic relation at arbitrary endpoints before any coefficient
is extracted.

The underlying geometric kernel and its normalized Mellin interpolation
are already present in Ihara, Nakamura, and Yamamoto
\cite[Proposition~3.3 and Theorem~4.1]{INY}; the published paper is
\emph{The Ramanujan Journal} \textbf{67} (2025), article~1,
doi:10.1007/s11139-025-01045-2.  In their notation the kernel below is
$e^{-b\sum_jx_j}G_k(e^{-x_1},\ldots,e^{-x_k};a-b)$.
Equivalently it is $e^{-(b-1)\sum_jx_j}$ times their $b=1$ kernel.
Its divided-difference form, composition proof, and continuation proof
are included to make the normalization and later estimates explicit;
the underlying interpolant and its geometric Mellin representation are
not claimed as new.  The developments here are the simultaneous twist
tube, the full dimension-independent zero-decoration polydisc, the
all-block identity with arbitrary independent free orders, and its
compatibility with the transverse derivative calculus developed below.

Write $\mathcal H_{a,b}(s_1,\ldots,s_k)$ for the entire relative
Hurwitz interpolant, normalized by
\[
 Z_b(w)=\sum_{w=uv} Z_a(u)\mathcal H_{a,b}(v),\qquad
 Z_c(s_1,\ldots,s_j)=\sum_{n_1>\cdots>n_j\geq0}
                         \prod_{\ell=1}^j(n_\ell+c)^{-s_\ell}.
\]
Empty words have value one.  Throughout this section
$\Re a,\Re b>0$, and powers of the positive-real-axis variables use
their real logarithms.  All endpoint powers use the branches obtained
on the right half-plane.

\subsection{The established geometric kernel in divided-difference form}

For $\boldsymbol x=(x_1,\ldots,x_k)$ put
\[
 X_0=0,\qquad X_j=x_1+\cdots+x_j,\qquad Q_j=e^{-X_j},\qquad
 \Delta=a-b.
\]
At pairwise distinct nodes define
\begin{equation}\label{zbg:kernel}
 \mathscr K_{a,b}^{(k)}(\boldsymbol x)
 =e^{-bX_k-\sum_{j=1}^{k-1}X_j}
       \sum_{j=0}^k
       \frac{e^{-\Delta X_j}}{\prod_{\ell\ne j}(Q_j-Q_\ell)}.
\end{equation}
The sum is the divided difference
$[Q_0,\ldots,Q_k]z^\Delta$, where the value at $Q_j$ is specified
as $e^{-\Delta X_j}$.  This specification keeps track of the
logarithms even when the exponential nodes wind around zero.  Set
$\mathscr K^{(0)}=1$.

\begin{lemma}[Kernel analyticity and exponential control]\label{zbg:bounds}
The expression \eqref{zbg:kernel} extends holomorphically through
every hyperplane $X_j-X_i=0$, away from the nonzero-period
hyperplanes below. Its only possible finite singular
hyperplanes are
\[
 X_j-X_i=2\pi i\ell,\qquad 0\leq i<j\leq k,
                   \quad \ell\in\mathbb Z\setminus\{0\}.
\]
Define the open convex tube
\begin{equation}\label{zbg:tube}
 \Omega_k=\left\{\boldsymbol t\in\mathbb C^k:
   \left|\Im\sum_{j=p}^q t_j\right|<2\pi
                      \quad(1\leq p\leq q\leq k)\right\}.
\end{equation}
On compact subsets of $\Re a,\Re b>0$ and
$\boldsymbol t\in\Omega_k$, for every multi-index $\boldsymbol\nu$
there are $C,\delta>0$ such that
\begin{equation}\label{zbg:bound}
 \left|\partial_{\boldsymbol x}^{\boldsymbol\nu}
       \mathscr K_{a,b}^{(k)}(\boldsymbol x-\boldsymbol t)\right|
       \leq C e^{-\delta(x_1+\cdots+x_k)},\qquad x_j\geq0.
\end{equation}
\end{lemma}

\begin{proof}
If nodes collide with $X_j-X_i=0$, their assigned logarithms agree.
On a sufficiently small neighborhood of each cluster, choose the common
holomorphic logarithm.  The ordinary confluent divided difference of
its holomorphic power function extends across all coinciding nodes.
Different clusters have nonzero denominator differences.  This proves
the simultaneous removability assertion, including intersections of
arbitrarily many zero hyperplanes away from the remaining singular
hyperplanes. Nonzero exponential periods are
exactly the remaining possible denominator zeros.

The first-order kernel is
\[
 \mathscr K_{a,b}^{(1)}(v)=\frac{e^{-bv}-e^{-av}}{1-e^{-v}}.
\]
The elementary divided-difference recurrence gives, for $k\geq2$,
\begin{equation}\label{zbg:recurrence}
 \mathscr K_{a,b}^{(k)}(v,w,\boldsymbol z)=
 \frac{e^{-v}\mathscr K_{a,b}^{(k-1)}(v+w,\boldsymbol z)
       -e^{-av}\mathscr K_{a,b}^{(k-1)}(w,\boldsymbol z)}
      {1-e^{-v}}.
\end{equation}
Indeed the two divided differences in the numerator omit $Q_1$ and
$Q_0$, respectively; scaling the nodes of the second one by $Q_1$
gives exactly the displayed exponential factor.

This recurrence proves \eqref{zbg:bound} by induction.  For large
$x_1$, the denominator is uniformly bounded away from zero, and both
numerator terms have the required exponential decay.  For bounded
$x_1$, the only possible denominator zero is the removable zero at
$v=0$; the numerator vanishes there because its two kernel terms agree.  A fixed small neighborhood of the shifted positive orthant
stays away from all nonzero-period hyperplanes, because the compact
subset in \eqref{zbg:tube} has a positive distance from its boundary.
The maximum principle, or the Cauchy formula around the removable
zero, then bounds the quotient by the inductively bounded numerator.
Merging $v+w$ preserves the tube inequalities, since they involve
consecutive sums.  One may take $\delta$ smaller than the minimum of
$1,\Re a,\Re b$ on the endpoint compact set.  Cauchy estimates in the
same uniform neighborhood prove the bounds for every derivative.
\end{proof}

\subsection{Why the kernel gives the specified noninteger interpolant}

Let
\begin{equation}\label{zbg:infinite}
 \mathscr B_c^{(k)}(\boldsymbol x)
 =\frac{e^{-cX_k-\sum_{j=1}^{k-1}X_j}}
              {\prod_{j=1}^k(1-e^{-X_j})},\qquad x_j>0.
\end{equation}
This is the geometric kernel of the ordered infinite sum $Z_c$:
introducing the positive successive gaps gives
$\mathscr B_c^{(k)}=\sum_{n_1>\cdots>n_k\geq0}
e^{-\sum_j(n_j+c)x_j}$.

\begin{lemma}[Exact endpoint composition]\label{zbg:composition}
For arbitrary right-half-plane endpoints $a,b,c$,
\begin{equation}\label{zbg:kernelcomposition}
 \mathscr K_{a,b}^{(k)}(x_1,\ldots,x_k)
 =\sum_{j=0}^k
   \mathscr K_{a,c}^{(j)}(x_1,\ldots,x_j)
   \mathscr K_{c,b}^{(k-j)}(x_{j+1},\ldots,x_k).
\end{equation}
Consequently,
\begin{equation}\label{zbg:convolution}
 \mathscr B_b^{(k)}=\sum_{j=0}^k
                       \mathscr B_a^{(j)}\mathscr K_{a,b}^{(k-j)},
\end{equation}
with the same prefix/suffix arguments.
\end{lemma}

\begin{proof}
Apply the divided-difference product identity
\[
 [Q_0,\ldots,Q_k](fg)=\sum_{j=0}^k
 [Q_0,\ldots,Q_j]f\,[Q_j,\ldots,Q_k]g
\]
to $f(z)=z^{a-c}$ and $g(z)=z^{c-b}$.  The suffix nodes in its
kernel are $Q_\ell/Q_j$.  Scaling a divided difference of order
$k-j$ supplies the factor $Q_j^{k-j-(c-b)}$; together with the
kernel prefactors this makes every summand have the common factor
$e^{-bX_k-\sum_{\ell=1}^{k-1}X_\ell}$.  Thus the product identity
is exactly \eqref{zbg:kernelcomposition}.  This argument is first
made at distinct positive nodes and then continued through the
removable collisions.

For $x_j>0$, let a real upper endpoint $A$ tend to $+\infty$ in
$\mathscr K_{A,b}=\mathscr K_{A,a}*\mathscr K_{a,b}$.  In
\eqref{zbg:kernel}, every term except $j=0$ tends to zero.  The
remaining term is precisely \eqref{zbg:infinite}.  This gives
\eqref{zbg:convolution}, a finite identity requiring no interchange
with an infinite depth sum.  In particular, the proof is valid for
noninteger $a-b$ and does not use uniqueness of interpolation from
integer values.
\end{proof}

\subsection{The entire multiple-Lerch completion}

For $\Re s_j>0$ define
\begin{equation}\label{zbg:Mellin}
 \mathscr E_{a,b}(\boldsymbol s;\boldsymbol t)
 =\frac1{\prod_{j=1}^k\Gamma(s_j)}
   \int_{(0,\infty)^k}\prod_{j=1}^k x_j^{s_j-1}
      \mathscr K_{a,b}^{(k)}(\boldsymbol x-\boldsymbol t)
           \,d\boldsymbol x,\qquad\boldsymbol t\in\Omega_k.
\end{equation}

\begin{theorem}[Independent spectral orders and simultaneous branch cancellation]
\label{zbg:entire}
The function \eqref{zbg:Mellin} extends jointly to an entire function
of $\boldsymbol s\in\mathbb C^k$, holomorphic in
$\boldsymbol t\in\Omega_k$, with locally normal derivatives of
every fixed spectral and decoration order.  It satisfies
\begin{equation}\label{zbg:lowering}
 \mathscr E_{a,b}(\boldsymbol s;\boldsymbol0)
       =\mathcal H_{a,b}(\boldsymbol s),\qquad
 \partial_{t_j}\mathscr E_{a,b}(\boldsymbol s;\boldsymbol t)
       =\mathscr E_{a,b}(\boldsymbol s-\boldsymbol e_j;\boldsymbol t).
\end{equation}
In a neighborhood of $\boldsymbol t=0$ its normally convergent Taylor
series is
\begin{equation}\label{zbg:Taylor}
 \mathscr E_{a,b}(\boldsymbol s;\boldsymbol t)
 =\sum_{\boldsymbol\nu\geq0}
     \mathcal H_{a,b}(\boldsymbol s-\boldsymbol\nu)
                          \frac{\boldsymbol t^{\boldsymbol\nu}}{\boldsymbol\nu!}.
\end{equation}
For instance, the Taylor series converges normally whenever
$\sum_j|t_j|<2\pi$.
\end{theorem}

\begin{proof}
Lemma~\ref{zbg:bounds} proves absolute convergence of the integral.
Repeated integration by parts gives the particularly useful continuation
formula
\begin{equation}\label{zbg:continuation}
 \mathscr E_{a,b}(\boldsymbol s;\boldsymbol t)
 =\frac{(-1)^{|\boldsymbol N|}}
             {\prod_j\Gamma(s_j+N_j)}
    \int_{(0,\infty)^k}\prod_jx_j^{s_j+N_j-1}
      \partial_{\boldsymbol x}^{\boldsymbol N}
          \mathscr K_{a,b}^{(k)}(\boldsymbol x-\boldsymbol t)
                                                    \,d\boldsymbol x,
\end{equation}
valid for $\Re s_j>-N_j$.  It agrees with the original integral
where $\Re s_j>0$; the endpoint terms there vanish.  Since each
$N_j$ is arbitrary, these compatible formulas prove joint entireness.
The same exponential majorants, with logarithmic powers from spectral
differentiation, prove local normality.

At $\boldsymbol t=0$, apply the normalized Mellin transform to
\eqref{zbg:convolution}, initially with all spectral real parts
sufficiently large.  The product of prefix/suffix transforms is a
product of independent integrals.  The result is exactly the defining
deconcatenation relation for $\mathcal H_{a,b}$.  This proves the
first identity in \eqref{zbg:lowering}; entireness extends it everywhere.
The equality $\partial_{t_j}\mathscr K=-\partial_{x_j}\mathscr K$
and one integration by parts prove the lowering relation first where
$\Re s_j>1$, and then everywhere.  Taylor's theorem proves
\eqref{zbg:Taylor}.  Every polydisc whose coordinate radii have sum
less than $2\pi$ lies in $\Omega_k$, giving the stated domain.
\end{proof}

For $\boldsymbol t\in\Omega_k$ with $\Re t_j<0$ for every $j$, let
\[
 \mathscr Z_c(\boldsymbol s;\boldsymbol t)
 =\sum_{n_1>\cdots>n_k\geq0}
       \prod_j(n_j+c)^{-s_j}e^{(n_j+c)t_j}.
\]
These ordered multiple-Lerch series (with the standard depth-one
normalization \cite[\S25.14]{DLMFLerch}) are absolutely convergent and
entire in their independent orders.  In that domain,
$\mathscr E_{a,b}=\mathscr Z_a^{-1}*\mathscr Z_b$ under
deconcatenation, by the same kernel identity.  At each fixed depth the
inverse is a finite expression.  For example,
\begin{align*}
 \mathscr E_{a,b}(s,t;u,v)={}&
  \mathscr Z_b(s,t;u,v)-\mathscr Z_a(s,t;u,v)\\
 &-\mathscr Z_a(s;u)
               \bigl(\mathscr Z_b(t;v)-\mathscr Z_a(t;v)\bigr).
\end{align*}
Theorem~\ref{zbg:entire} proves that the whole expression has a
single holomorphic continuation throughout $\Omega_k$, including all
simultaneous zero-twist branch crossings.  Continuing the summands
independently without this completion is unnecessary.

\subsection{Summing every zero block on a dimension-independent polydisc}

Let $\boldsymbol u=(u_0,\ldots,u_k)$, $|u_j|<1$, and take the
logarithms $L_j=\log(1+u_j)$ determined at the origin.  Put
$t_j=L_j-L_{j-1}$.

\begin{theorem}[All zero decorations around independent orders]
\label{zbg:main}
For every $k\geq1$,
\begin{align}\label{zbg:generator}
 &\sum_{p_0,\ldots,p_k\geq0}
 \mathcal H_{a,b}(\{0\}^{p_0},s_1,\{0\}^{p_1},\ldots,
                       s_k,\{0\}^{p_k})\prod_{j=0}^ku_j^{p_j}
 \notag\\
 &\qquad=(1+u_0)^{a-1}(1+u_k)^{-b}
                    \prod_{j=1}^{k-1}(1+u_j)^{-1}
        \mathscr E_{a,b}(\boldsymbol s;
                L_1-L_0,\ldots,L_k-L_{k-1}).
\end{align}
The series converges absolutely and locally normally on the full
polydisc $|u_j|<1$, uniformly on compact spectral and endpoint sets.
Every fixed collection of independent spectral derivatives may be
taken term by term.
\end{theorem}

\begin{proof}
For $|u|<1$, $1+u$ lies in the right half-plane, so
$|\Im\log(1+u)|<\pi/2$.  Every consecutive sum of the $t_j$
telescopes to $L_q-L_{p-1}$ and has imaginary part of absolute value
less than $\pi$.  Thus the right side of \eqref{zbg:generator}
is holomorphic on the entire stated polydisc, independently of $k$.

To identify its Taylor coefficients, use \eqref{zbg:Taylor} locally
and formally replace
$\mathcal H_{a,b}(\boldsymbol s-\boldsymbol\nu)$ by
$\boldsymbol y^{\boldsymbol\nu}$.  The resulting expression is
\[
 (1+u_0)^{a-y_1-1}
 \prod_{j=1}^{k-1}(1+u_j)^{y_j-y_{j+1}-1}
 (1+u_k)^{y_k-b}.
\]
Its coefficient of $\prod u_j^{p_j}$ is
\[
 \binom{a-y_1-1}{p_0}
 \prod_{j=1}^{k-1}\binom{y_j-y_{j+1}-1}{p_j}
 \binom{y_k-b}{p_k}.
\]
This is exactly the canonical strict-gap polynomial of the zero
blocks.  The analytic gap-deletion theorem in \cite{IncomingParity},
with the independent-order continuation fixed in \cite{IncomingIndependent},
identifies its finite polynomial expansion with
the left coefficient in \eqref{zbg:generator}.  Each coefficient
uses only finitely many terms of \eqref{zbg:Taylor}; no interchange
of a divergent formal order-shift sum is involved.  Holomorphy of the
right side now proves absolute local convergence of its Taylor series
throughout the unit polydisc.  Spectral derivatives obey the same
Cauchy estimates uniformly on compact spectral sets.
\end{proof}

\begin{corollary}[Insertion sums and normalized primitives]\label{zbg:corollaries}
For every integer $r\geq0$,
\begin{equation}\label{zbg:insertion}
 \sum_{p_0+\cdots+p_k=r}
 \mathcal H_{a,b}(\{0\}^{p_0},s_1,\ldots,s_k,\{0\}^{p_k})
 =\binom{a-b-k}{r}\mathcal H_{a,b}(s_1,\ldots,s_k),
\end{equation}
where the zero blocks between consecutive displayed free orders are
included in the summand as in \eqref{zbg:generator}.  Moreover,
\begin{equation}\label{zbg:primitive}
 \int_{\tau_0}^{t_j}\mathscr E_{a,b}(\boldsymbol s;
                      t_1,\ldots,\tau,\ldots,t_k)\,d\tau
 =\mathscr E_{a,b}(\boldsymbol s+\boldsymbol e_j;\boldsymbol t)
  -\mathscr E_{a,b}(\boldsymbol s+\boldsymbol e_j;
                     t_1,\ldots,\tau_0,\ldots,t_k),
\end{equation}
along any path in the indicated slice of $\Omega_k$.
For $\boldsymbol m\in\mathbb N_0^k$,
\begin{equation}\label{zbg:negative}
 \mathscr E_{a,b}(-\boldsymbol m;\boldsymbol t)
 =(-1)^{|\boldsymbol m|}
       \partial_{\boldsymbol x}^{\boldsymbol m}
            \mathscr K_{a,b}^{(k)}(-\boldsymbol t).
\end{equation}
\end{corollary}

\begin{proof}
Set every $u_j=u$ in \eqref{zbg:generator}.  All twists vanish and
the prefactor becomes $(1+u)^{a-b-k}$.  The primitive is
\eqref{zbg:lowering} with the additive normalization retained.
For \eqref{zbg:negative}, use \eqref{zbg:continuation} with
$N_j=m_j+1$ and integrate the derivatives successively.  Exponential
decay removes the upper boundary terms and supplies precisely one
additional minus sign for each integration.
\end{proof}

At $a=b+N$ with $N\geq0$ integral, the kernel is the finite ordered
exponential sum and $\mathscr E_{a,b}$ is its weighted finite
Dirichlet sum.  The generating function is consequently a polynomial
of total degree at most $N-k$ if $N\geq k$, and is zero if $N<k$.
For $k=1$ the theorem specializes to the previously proved
single-Lerch generator.  For $k>1$ it resolves the requested joint
convergence and simultaneous branch-cancellation problem.  It supplies
an explicit Gamma-separated kernel, with no claim that this kernel is
minimal in an unspecified class.
